\section{Background}

\subsection{Perceptual image hashes}

All four of the paper's hashes come from Buchner's \texttt{imagehash} library~\cite{imagehash},
which also produced the paper's reported numbers. Three of them are fixed-width binary
strings compared with the Hamming distance (the number of differing bits):
\begin{itemize}
  \item \textbf{aHash} (64 bits): shrink to $8\times8$ grayscale and set each bit by whether
        the pixel is brighter than the mean. A coarse brightness layout.
  \item \textbf{pHash} (64 bits): shrink to $32\times32$, take the discrete cosine transform,
        and threshold the $8\times8$ lowest frequencies at their median. Coarse shapes.
  \item \textbf{hsvHash} (42 bits): a global histogram of the image's colours, with black, gray
        and saturated pixels binned by hue. The paper's text describes a block-based HSV hash, but the
        \texttt{colorhash()} function that generated its numbers is global. We followed the
        code, since that is where the paper's numbers came from.
\end{itemize}

\textbf{sHash} is different. It segments the image into its largest regions and computes a
64-bit dHash (a horizontal-gradient hash) of each region's bounding box. One image therefore
yields a \emph{list} of hashes whose length depends on the image. The paper adds it
specifically to survive cropping: if part of the image is cut away, the remaining large
objects should still match.

\subsection{The paper's detector}

Each hash votes independently: it votes ``copy'' when the query's nearest stored image is
within that hash's threshold. The \emph{2-Minimal} detector reports a copy when at least two
of the four vote, which keeps false alarms low because no single hash has to be trusted on
its own (Figure~\ref{fig:paper}). Formally, with $d_i$ the distance under hash $i$ and
$t_i$ its threshold,
\[
  \mathrm{copy}(q,r) \;=\; \Bigl[\, \textstyle\sum_{i=1}^{4} \mathbb{1}\bigl[d_i(q,r)\le t_i\bigr] \;\ge\; 2 \,\Bigr].
\]
The thresholds given in the paper are aHash $\le 7$, pHash $\le 15$, hsvHash $\le 3$ and sHash
$\le 17$~\cite{kotzer2025}.

\begin{figure}[ht]
\centering
\begin{tikzpicture}[
  font=\small,
  box/.style={draw=ink2, rounded corners=2pt, minimum height=8mm, align=center, fill=white},
  sig/.style={box, minimum width=27mm},
  >={Stealth[length=2mm]}]
  \node[box, minimum width=22mm] (q) {Query image\\[-1pt]\footnotesize a new mint};
  \node[sig, right=14mm of q, yshift=16.5mm] (a) {aHash\\[-1pt]\footnotesize brightness layout};
  \node[sig, below=2mm of a] (p) {pHash\\[-1pt]\footnotesize coarse shapes};
  \node[sig, below=2mm of p] (h) {hsvHash\\[-1pt]\footnotesize colour mix};
  \node[sig, below=2mm of h, draw=shashorange, very thick] (s) {sHash\\[-1pt]\footnotesize image pieces (cropping)};
  \foreach \n in {a,p,h,s} {
    \node[box, right=5mm of \n, minimum width=17mm, fill=rule!50] (t\n) {BK-tree};
    \draw[->] (q.east) -- (\n.west);
    \draw[->] (\n) -- (t\n);
  }
  \node[box, right=13mm of tp, yshift=-5mm, minimum width=22mm, fill=technion!8] (v) {$\ge 2$ of 4\\ agree?};
  \foreach \n in {a,p,h,s} \draw[->] (t\n.east) -- (v.west);
  \node[box, right=7mm of v, minimum width=14mm] (out) {copy /\\ not a copy};
  \draw[->] (v) -- (out);
\end{tikzpicture}
\caption{The paper's 2-Minimal detector. Four hashes, each searched in its own BK-tree; a
copy needs two votes. Our project replaces the highlighted sHash.}
\label{fig:paper}
\end{figure}

\paragraph{Storage on the blockchain.} The paper proposes adding the four hashes as metadata
fields of every mint transaction, so a validator can check a new mint ``in a reasonable time''
without fetching images. It notes that a transaction is usually 300--500 bytes and that each
hash needs only a few bytes, so the overhead is low~\cite{kotzer2025}. We use 400 bytes as a
reference transaction. The four hashes take about 56 bytes: 24 fixed for aHash, pHash and
hsvHash together, plus a median of 32 for sHash's variable-length list
(Section~\ref{sec:storage}).

\subsection{BK-trees need a metric}

A BK-tree~\cite{burkhard1973} stores one item per node and labels each child edge with the
child's distance from its parent. To find everything within distance $t$ of a query $q$, the
search computes $d=d(q,\text{node})$ and descends only into child edges labelled in
$[d-t,\,d+t]$. This pruning is what makes the tree fast, and it is only safe when the
distance is a \emph{metric}: in particular it must be symmetric, $d(x,y)=d(y,x)$, and obey the
triangle inequality, $d(x,z)\le d(x,y)+d(y,z)$. The paper states this requirement itself.
Hamming distance on fixed-width strings is a metric. A distance between variable-length lists
is not automatically one.

\subsection{ORB, LSH and RANSAC}

ORB (Oriented FAST and Rotated BRIEF)~\cite{rublee2011} finds \emph{keypoints}, corners and
sharp junctions that can be found again after an edit, and describes the neighbourhood of
each with a 256-bit (32-byte) binary descriptor. Each keypoint also gets an orientation, which
is what lets matches survive rotation. Two images are compared by matching descriptors in
Hamming distance. Many matches are coincidental, so RANSAC~\cite{fischler1981} searches for
the single geometric transformation (here a homography) that the most matches agree on and
discards the rest. The number of agreeing matches, the \emph{inliers}, is the similarity
score.

Because one image yields hundreds of descriptors rather than one string, ORB cannot be put in
a BK-tree. The natural index for binary descriptors is locality-sensitive hashing
(LSH)~\cite{gionis1999}, which places similar descriptors in the same bucket so that only a
short list of candidates needs RANSAC verification.
